\scanpage{99}
\mspage{1}
\uuline{Théorème}\qquad $X/k$, $k$ corps \uuline{alg. clos}

$X$ \deleted{\uuline{intègre}}\inserted{\uline{normal} \uline{\unsure{irred|irré d|intègre}}} \uuline{projectif}, $X\subset\mathbb P^{r}_{k}$

\deleted{$\operatorname{prof}X\geq2$ (d'} $\dim X\geq2$ (d'où $\operatorname{prof}X\geq2$)

$G$ groupe algébrique \uuline{affine} \uuline{sans \unsure{composants|composantes}} $G_{m}$. (p.ex. un groupe alg. fini).

$K=\deleted{\unsure{\overline{\phantom{k(t_1,t_r)}}}}k(t_{1},\ldots,t_{r})$, $\overline K$ sa clôture algébrique

$Y_{\overline K}$ la section hyperplane générique, définie sur $\overline K$ :

\[
\text{Alors}\qquad H^{1}(X,G)\longrightarrow H^{1}(Y_{\overline K},G)\qquad\text{\uuline{est injectif}}
\]

La structure des groupes algébriques comm. nous ramène, aux cas simples $G_{a}$, \deleted{\unsure{$G_{m}$}} groupe fini séparable, \inserted{$\alpha_{p}$}, $\mu_{p}$. Donc on est ramené à prouver

% slanted
\margin{N.B. Généralisation au cas d'un morphisme donné $X\longrightarrow P^{r}$\ldots.}

\uline{Corollaire 1}\qquad $H^{1}(X,\underline{\mathcal O}_{X})\longrightarrow H^{1}(Y_{\overline K},\underline{\mathcal O}_{Y_{\overline K}})$ est injectif

\deleted{\uline{Corollaire 2}\qquad $\mathrm{Pic}(X)\longrightarrow\mathrm{Pic}(Y_{\overline K})$ est injectif}

\uline{Corollaire 2}\qquad $\pi_{1}(Y_{\overline K})\longrightarrow\pi_{1}(X)$ est surjectif [c'est \uuline{Bertini}]

\uline{Corollaire \deleted{\unsure{2|3}}3}\qquad $H^{1}(X,\alpha_{p})\unsure{-|\to}H^{1}(X_{\overline K},\alpha_{p})$ injectif

\uline{Corollaire \deleted{\unsure{3|4}}4}\qquad $H^{1}(X,\mu_{p})\unsure{-|\to}H^{1}(Y_{\overline K},\mu_{p})$ injectif.

On part du corollaire 2, connu grâce à Bertini, \uline{ou} Lefschetz-Grauert [utiliser \deleted{\unsure{disque|doit}} $\operatorname{prof}X\geq2$]. Les en déduit $H^{1}(X,\underline{\mathcal O}_{X})^{F}\Rightarrow H^{1}(X\deleted{\unsure{,}}\mathbb Z/p\mathbb Z)\longrightarrow H^{1}(Y_{\overline K},\underline{\mathcal O}_{Y_{\overline K}})^{F}$ est injectif, on est donc ramené pour le cor. 1 de prouver que $H^{1}(X,\underline{\mathcal O}_{X})^{1-F}-H^{1}(Y_{\overline K},\underline{\mathcal O}_{Y_{\overline K}})^{1-F}$ est injectif [N.B., en tous les \unsure{;} le noyau est stable sous $F$\ldots] i.e. au corollaire 3. Donc on est ramené \deleted{\unsure{au}}\inserted{[en car. $p>0$]} \deleted{\unsure{à}} au cas des cor. 3. et 4. En car $p=0$,

% archive sheet number: 98

